\documentclass[letterpaper]{article}
\usepackage[utf8x]{inputenc}
\usepackage[T1]{fontenc}
\usepackage[activeacute,english]{babel}
\usepackage[vmargin=4cm,tmargin=3cm,hmargin=2cm,letterpaper]{geometry}%
\usepackage{helvet}
\usepackage{amsmath,amsfonts,amssymb}
\usepackage{graphicx}
\usepackage{color}
\usepackage{xcolor}
\usepackage{verbatim}
\usepackage{tabls}
\usepackage{lastpage}
\usepackage{fancyhdr}
\usepackage{url}
\usepackage{listings}
%%%%%%%%%%%%%%%%%%%%%%%%%%%%%%%%%%%%%%%%%%%%%%%%%%%%%%%%%%%%%%%%%%%%%%%%%%%%%%%%%%%%%%%
\usepackage{tikz}
\usepackage{pgf}
\usepackage{pgffor}
\usepgfmodule{plot}
\usepackage{wrapfig}
\usetikzlibrary{arrows,decorations,snakes,backgrounds,fit,calc,through,scopes,positioning,automata,chains,er,fadings,calendar,matrix,mindmap,folding,patterns,petri,plothandlers,plotmarks,shadows,shapes,shapes.arrows,topaths,trees}

\lstset{% general command to set parameter(s)
%   basicstyle=\small,
  % print whole listing small
%   keywordstyle=\color{black}\bfseries\underbar,
  % underlined bold black keywords
%   identifierstyle=,
  % nothing happens
%   commentstyle=\color{white}, % white comments
%   stringstyle=\ttfamily,
  % typewriter type for strings
  showstringspaces=false}
  % no special string spaces
%\DeclareGraphicsExtensions{.bmp,.png,.pdf,.jpg,.eps}
\pagestyle{fancy}
\color{black}
\fancyhead{}
\renewcommand{\headrule}{\hrule\vspace*{0.5mm}\rule{\linewidth}{0.8mm}}
\renewcommand{\familydefault}{\sfdefault}

\graphicspath{{/home/luis/Documents/estructuras}}
\lhead{\includegraphics[width=2cm]{logoucr.eps}}
\rhead{\includegraphics[width=3cm]{eie-text-gray-6x3cm.eps}}
\chead{UNIVERSITY OF COSTA RICA\\FACULTY ENGINEERING\\ELECTRICAL ENGINENEERING SCHOOL\\\textbf{DATA STRUCTURES AND ALGORITHMS ABSTRACT\\ FOR ENGINEERING}\\IE-0217\\CYCLE I 2014\\BIBLIOGRAPHIC RESEARCH 1}

\lfoot{}%
\cfoot{}%
%\cfoot{\thepage\ de \pageref{LastPage}}%
\rfoot{}%

%%%%%%%%%%%%%%%%%%%%%%%%%%%%%%%%%%%%%%%%%%%%%%%%%%%%%%%%%%%%%%%%%%%%%%%%%%%%%%%%%%%%%%%%%%%%%%%%%%%%%%%%%%%%%%%
\newcommand{\uic}{black} %user-input color
%%%%%%%%%%%%%%%%%%%%%%%%%%%%%%%%%%%%%%%%%%%%%%%%%%%%%%%%%%%%%%%%%%%%%%%%%%%%%%%%%%%%%%%%%%%%%%%%%%%%%%%%%%%%%%%%%%
\newcommand{\uim}{\\} %user-input marker
%%%%%%%%%%%%%%%%%%%%%%%%%%%%%%%%%%%%%%%%%%%%%%%%%%%%%%%%%%%%%%%%%%%%%%%%%%%%%%%%%%%%%%%%%%%%%%%%%%%%%%%%%%%%%%%%%%
\newcommand{\userinput}[1]{\textcolor{\uic}{\uim#1\uim}}

%%%%%%%%%%%%%%%%%%%%%%%%%%%%%%%%%%%%%%%%%%%%%%%%%%%%%%%%%%%%%%%%%%%%%%%%%%%%%%%%%%%%%%%%%%%%%%%%%%%%%%%%%%%%%%%%%%
\begin{document}\vspace*{2cm}
%%%%%%%%%%%%%%%%%%%%%%%%%%%%%%%%%%%%%%%%%%%%%%%%%%%%%%%%%%%%%%%%%%%%%%%%%%%%%%%%%%%%%%%%%%%%%%%%%%%%%%%%%%%%%%%%%%

%%%%%%%%%%%%%%%%%%%%%%%%%%%%%%%%%%%%%%%%%%%%%%%%%%%%%%%%%%%%%%%%%%%%%%%%%%%%%%%%%%%%%%%%%%%%%%%%%%%%%%%%%%%%%%%%%%
\begin{center}
\Huge
\userinput{Genetic Algorithms}
\vspace*{1cm}
\end{center}

\noindent
\small\baselineskip=14pt
\textbf{Student:} \userinput{Ernesto Céspedes Montero, Ricardo Chacón Carvajal, Luis Sergio Yannicelli Vargas}\\
%\textbf{License:} \userinput{A87086}\\
\textbf{Email:} \textit{netosoy@gmail.com, vrlrdo@gmail.com, lsyannnicelli@gmail.com}\\

\section{Introduction}
In the present days, technologic growth specially in the programming field, represents and important issue, taking big part on solving a variety of topics, implementing algorithms as a tool. * Genetic algorithms are evolutionary algorithms, more broadly belonging to the Biologically Inspired Computation, useful for optimization problems. Its structure is inspired firmly on the field of genetics and Mendelian biology, specifically populations, gene frequency, chromosomes, genes, alleles, and so on in terms related to genetics; on the other hand it applies genetic processes such as evolution, recombination, mutation, heredity and others to improve the algorithm's efficiency for its goal which is combinatorial optimization.

\section{Objectives}

\subsection{General Objective}

To implement a C++ library of genetic algorithms.

\subsection{Specific Objectives}

Specific Objectives are:\\

\begin{enumerate}
\item To understand genetics algorithms method as evolutive algorithm.
\item To describe genetic algorithms method and some aplications in which this algorithm could be efficient.
\item To implement an C++ library of genetic algorithms.
\end{enumerate}

\section{Methodology}

La finalidad de proyecto será la de crear una librería abierta de algoritmos genéticos en lenguaje C++, en otras palabras, se pretende crear un directorio de archivos de cabecera (.h) que contenga todas las declaraciones de las funciones de la librería y otro directorio de archivos fuente (.cpp) que contenga la implementación de las funciones. La librearía a crear será una librería dinámica por lo que se pretende crear un archivo .so final, el cual será compilado por medio de la herramienta CMake.

La implementación de las funciones se basará en estudios bibliográficos sobre la teoría general de los algoritmos genéticos para entender la estructura de dato a implementar en las funciones de inicialización, evaluación, selección, cruzamiento, mutuación y remplazo. Además se pretende buscar pseudocódigos o implementaciones en otros lenguajes de las funciones mencionadas, para que junto con la investigación bibliográfica, se comprenda la estructura de estas y se esté en capacidad de implementar de manera de manera propia los algoritmos mencionados.



%\begin{figure}
%\centering
%\includegraphics[width=0.8\linewidth]{ejemplo_grafo2.eps}
%\caption{Grapho's Example}
%\label{F:ejemplo_grafo}
%\end{figure}

\begin{thebibliography}{99}
\bibitem{Floyd}
Consultado el 25 de marzo de 2014, en: http://dalila.sip.ucm.es/~manuel/Informatica/FloydWarshall.pdf

\bibitem{Grafo}
Consultado el 25 de marzo de 2014, en: http://www.fing.edu.uy/inco/grupos/bioinf/bioinfo1/teorico/grafos.pdf

\end{thebibliography}

\end{document}
